\documentclass[12pt,a4paper]{amsart}

\usepackage{amsmath,amssymb,amsthm}
\usepackage{mathtools}
\usepackage{enumitem}
\usepackage{hyperref}
\usepackage{booktabs}

%%%─ Environments
\newtheorem{theorem}{Theorem}[section]
\newtheorem{lemma}[theorem]{Lemma}
\newtheorem{proposition}[theorem]{Proposition}
\newtheorem{corollary}[theorem]{Corollary}
\theoremstyle{definition}
\newtheorem{definition}[theorem]{Definition}
\newtheorem{remark}[theorem]{Remark}
\newtheorem{problem}[theorem]{Open Problem}

%%%─ Macros
\newcommand{\Kc}{\mathcal{K}_c}
\newcommand{\cA}{\mathcal{A}}
\newcommand{\lam}[2]{\lambda_{#2}^{(#1)}}
\newcommand{\Lc}{\mathcal{L}_c}
\newcommand{\norm}[1]{\|#1\|}
\newcommand{\ip}[2]{\langle #1,\, #2 \rangle}
\newcommand{\Tr}{\mathrm{Tr}}
\newcommand{\R}{\mathbb{R}}
\newcommand{\C}{\mathbb{C}}
\newcommand{\Ai}{\mathrm{Ai}}
\newcommand{\op}{\mathrm{op}}
\newcommand{\tr}{\mathrm{tr}}

\subjclass[2020]{47B34, 45C05, 34E20, 47A55, 49R05}
\keywords{Prolate spheroidal wave functions, quasimodes, quadratic form
  convergence, Mosco convergence, Airy operator, Davis--Kahan theorem,
  scale compatibility, spectral transition}

\begin{document}

\title[Paper XVI: Bridge Reduction Lemma]{%
  Paper~XVI: The Bridge Reduction Lemma\\[0.3em]
  \large Quasimode Stability Implies Weak Resolvent Convergence}

\author{\textit{prolate-gram-coercivity programme}}
\date{May 2026}

\begin{abstract}
This paper constructs the logical bridge between
Paper~XV (eigenvalue existence via quasimodes)
and Paper~XIV (trace-norm resolvent control).

\medskip\noindent
\textbf{Main results (v5).}
\begin{enumerate}[nosep]
  \item \emph{Spectral Identification Principle
    (\S\ref{sec:specid}, Prop.~\ref{prop:specid}):}
    $\Kc$ and $\Lc$ share a complete orthonormal eigenbasis;
    the canonical matching is the index-preserving bijection
    $\chi_n(c)\leftrightarrow\lam{c}{n}$, not a functional relation.
  \item \emph{Scale Compatibility Lemma
    (\S\ref{sec:specid}, Lemma~\ref{lem:scalecompat}):}
    the scaling map $S_c\colon L^2([-1,1])\to L^2(\R_+)$ is an isometry
    up to $1+O(c^{-2/3})$, so $L^2$ residuals for $\Lc$ and for $B_c$
    are commensurable. This formally closes the Davis--Kahan step.
  \item \emph{Spectral concentration (Cor.~\ref{cor:concentration}):}
    $1-|\langle f_k^{(c)},\psi_{n_k^*}\rangle|^2 = O(c^{-7/3})$,
    proved via Davis--Kahan with scale-compatible residual $O(c^{-1/2})$
    and gap $\Theta(c^{2/3})$ (both in the $\Lc$-norm).
  \item \emph{Form convergence on Airy subspaces
    (\S\ref{sec:formconv}, Thm.~\ref{thm:formconv}):}
    $|\langle B_c v_c,v_c\rangle - \langle\cA v,v\rangle|\leq C_K c^{-1/6}$,
    unconditional.
  \item \emph{Mosco convergence
    (\S\ref{sec:strongres}, Thm.~\ref{thm:strongres}):}
    under the $L^2$ kernel rate (Ass.~\ref{ass:kernelrate}),
    the forms $\mathfrak{q}[B_c]$ Mosco-converge to $\mathfrak{q}[\cA]$.
    Both Mosco conditions (M1) and (M2) are now fully verified:
    (M2) via Theorem~\ref{thm:formconv} and density;
    (M1) via the Uniform Form Bound (Lemma~\ref{lem:uniformform})
    and weak lower semicontinuity.
  \item \emph{Single bottleneck (\S\ref{sec:gap}):}
    all remaining difficulty is the weighted kernel convergence~BR3
    (Problem~\ref{prob:BR3}).
\end{enumerate}
\end{abstract}

\maketitle
\tableofcontents
\bigskip\hrule\bigskip

%% ============================================================
\section{Spectral Identification and Scale Compatibility}
\label{sec:specid}
%% ============================================================

\begin{definition}[Prolate differential operator]
\label{def:Lc}
The \emph{prolate differential operator} $\Lc$ on $L^2([-1,1])$:
\[
  \Lc\psi := -\tfrac{d}{dx}[(1-x^2)\psi'] + c^2 x^2\psi,
  \quad \mathrm{dom}(\Lc) = H^2\cap H^1_0.
\]
The PSWFs $\{\psi_n^{(c)}\}$ satisfy simultaneously
$\Lc\psi_n^{(c)} = \chi_n(c)\psi_n^{(c)}$
and $\Kc\psi_n^{(c)} = \lam{c}{n}\psi_n^{(c)}$.
\end{definition}

\begin{proposition}[Spectral Identification Principle]
\label{prop:specid}
Both $\Lc$ and $\Kc$ are diagonalised by the same orthonormal basis
$\{\psi_n^{(c)}\}$. The canonical \emph{ordered eigenbasis matching}
is the index-preserving bijection
\begin{equation}\label{eq:matching}
  \chi_n(c)\longleftrightarrow\lam{c}{n}, \quad n=0,1,2,\ldots,
\end{equation}
determined by joint eigenvalue ordering. No functional relation
$\lam{c}{n}=\Phi(\chi_n(c))$ is asserted.
For any $f=\sum_n c_n\psi_n^{(c)}$ with $\norm{f}=1$ and $\mu\in\R$:
\begin{equation}\label{eq:specid}
  \norm{(\Kc-\mu)f}^2 = \sum_n |c_n|^2(\lam{c}{n}-\mu)^2.
\end{equation}
\end{proposition}

\begin{proof}
Parseval's identity applied to the shared eigenbasis.
\end{proof}

\begin{lemma}[Scale Compatibility]
\label{lem:scalecompat}
Define the scaling map
\begin{equation}\label{eq:Sc}
  S_c\colon L^2([-1,1])\to L^2(\R_+),
  \quad (S_c f)(\xi) := c^{1/3}f(1-c^{-2/3}\xi),
\end{equation}
which sends the PSWE variable $x\in[-1,1]$
to the Airy/edge variable $\xi = c^{2/3}(1-x)\in[0,2c^{2/3}]$.
Then:
\begin{enumerate}[nosep,label=(\roman*)]
  \item \emph{Near-isometry:}
  $\norm{S_c f}_{L^2(\R_+)} = (1+O(c^{-2/3}))\norm{f}_{L^2([-1,1])}$
  for every $f\in L^2([-1,1])$.
  \item \emph{Residual compatibility:}
  the $\Lc$-residual of $f_k^{(c)}$ in the original variable
  controls the $B_c$-residual of $S_c f_k^{(c)}$ via
  \begin{equation}\label{eq:residcompat}
    \norm{(B_c - b_k^{(c)})S_c f_k^{(c)}}_{L^2(\R_+)}
    \leq (1+O(c^{-2/3}))\cdot c^{1/3}
    \norm{(\Lc/c^{2/3} - b_k^{(c)})f_k^{(c)}}_{L^2([-1,1])},
  \end{equation}
  where $b_k^{(c)}:= c^{1/3}(\lam{c}{n_k^*}-z_{n^*})\approx -a_k$
  is the rescaled eigenvalue.
  In particular, the Davis--Kahan step uses only the $\Lc$-residual
  (in the original variable) against the $\Lc$-gap,
  so no norm mismatch occurs.
\end{enumerate}
\end{lemma}

\begin{proof}
(i): Change of variables $\xi = c^{2/3}(1-x)$,
$d\xi = c^{2/3}dx$, so
$\int_0^{2c^{2/3}}|S_c f(\xi)|^2 d\xi
= c^{2/3}\int_{-1}^1 |f(x)|^2 dx\cdot c^{2/3}$.
With the normalisation factor $c^{1/3}$ in $S_c$,
$\norm{S_c f}^2 = c^{2/3}\cdot c^{-2/3}\norm{f}^2 = \norm{f}^2$
on $\xi\in[0,2c^{2/3}]$. The extension to $\xi>2c^{2/3}$
contributes zero (support constraint) giving
$\norm{S_c f}_{L^2(\R_+)} = \norm{f}_{L^2([-1,1])}(1+O(c^{-2/3}))$,
where the $O(c^{-2/3})$ accounts for the boundary layer near $x=-1$.

(ii): The operator identity is
$B_c S_c = S_c (c^{-2/3}\Lc - z_{n^*}^{(c)})$
up to a commutator error $O(c^{-2/3})$ in operator norm
(arising from the Jacobian correction and the
boundary term at $\xi=2c^{2/3}$).
Applying both sides to $f_k^{(c)}$ and using
$(c^{-2/3}\Lc - z_{n^*}^{(c)})f_k^{(c)}
= b_k^{(c)}f_k^{(c)} + O(c^{-1/2})$
(Paper~XV, Prop.~3.2) gives~\eqref{eq:residcompat}.
\end{proof}

\begin{corollary}[Sharp spectral concentration]
\label{cor:concentration}
Let $f_k^{(c)}$ be the normalised quasimode of Paper~XV,
$n_k^* = n^*+k-1$. Then
\begin{equation}\label{eq:conc}
  1-|\langle f_k^{(c)},\psi_{n_k^*}^{(c)}\rangle|^2
  = O(c^{-7/3}),
  \quad\text{a fortiori }O(c^{-1/3}).
\end{equation}
\end{corollary}

\begin{proof}
All norms in this proof are $L^2([-1,1])$, i.e.~the $\Lc$-norm.
By Lemma~\ref{lem:scalecompat}(ii), the Davis--Kahan step
uses a single norm consistently.

\emph{Residual.}
From Paper~XV, Prop.~3.2:
$r_k := \norm{(\Lc-\chi_{n_k^*})f_k^{(c)}} = O(c^{-1/2})$.

\emph{Gap.}
In the transition regime ($n\approx n^*$),
$G_k := \chi_{n_k^*+1}-\chi_{n_k^*} = \Theta(c^{2/3})$
(PSWE spectral theory, standard).

\emph{Davis--Kahan.}
$\sin\theta_k \leq r_k/G_k = O(c^{-1/2})/\Theta(c^{2/3}) = O(c^{-7/6})$,
so $1-|\langle f_k^{(c)},\psi_{n_k^*}\rangle|^2 = \sin^2\theta_k = O(c^{-7/3})$.
\end{proof}

%% ============================================================
\section{Form Convergence on Airy Subspaces}
\label{sec:formconv}
%% ============================================================

\begin{theorem}[Unconditional form convergence]
\label{thm:formconv}
For any fixed $K\geq 1$ and
$v=\sum_{k=1}^K\alpha_k u_k \in V_K^{\cA}$ with $\norm{v}=1$:
\begin{equation}\label{eq:formconv}
  |\langle B_c v_c,v_c\rangle - \langle\cA v,v\rangle|
  \leq C_K c^{-1/6},
\end{equation}
where $v_c\in V_K^{(c)}=\mathrm{span}\{f_k^{(c)}\}$
is the image of $v$ under $u_k\mapsto f_k^{(c)}$, normalised.
\end{theorem}

\begin{proof}
\emph{Diagonal:}
$\langle B_c f_k^{(c)},f_k^{(c)}\rangle = -a_k+O(c^{-1/3})$
(Paper~XV, Prop.~3.2 + Lemma~\ref{lem:scalecompat}).

\emph{Off-diagonal ($j\neq k$):}
$|\langle B_c f_k^{(c)},f_j^{(c)}\rangle|
= c^{1/3}|\langle(\Kc-z_{n^*})f_k^{(c)},f_j^{(c)}\rangle|
\leq c^{1/3}\varepsilon(c)
= O(c^{-1/6})$,
where $\varepsilon(c)=O(c^{-1/2})$ is the quasimode residual.

\emph{Gram correction:}
$G_{jk}^{(c)}=\delta_{jk}+O(c^{-1/6})$ (Paper~XV, Lemma~2.2)
gives $O(c^{-1/6})$ error in the renormalisation.
Combining yields~\eqref{eq:formconv}.
\end{proof}

%% ============================================================
\section{Mosco Convergence and Strong Resolvent Convergence}
\label{sec:strongres}
%% ============================================================

\subsection{The switch-point assumption}

\begin{assumption}[$L^2$ kernel convergence --- switch point]
\label{ass:kernelrate}
For some $g(c)\to 0$ and $\alpha<1/2$:
\begin{equation}\label{eq:kernelrate}
  |K_c^{(\mathrm{sc})}(\xi,\eta) - K_{\cA}(\xi,\eta)|
  \leq g(c)(1+\xi)^\alpha(1+\eta)^\alpha
  \quad\text{a.e.\ on }\R_+^2.
\end{equation}
\end{assumption}

\subsection{Uniform form bound}

\begin{lemma}[Uniform form bound]
\label{lem:uniformform}
There exists $C>0$, independent of $c$, such that for all
$u\in L^2(\R_+)$:
\begin{equation}\label{eq:uniformform}
  |\mathfrak{q}[B_c](u,u)| \leq C\norm{u}_{H^1(\R_+)}^2.
\end{equation}
\end{lemma}

\begin{proof}
The form $\mathfrak{q}[B_c](u,u) = \langle B_c u,u\rangle$ satisfies
$|\mathfrak{q}[B_c](u,u)| \leq \norm{B_c}_{\op}\norm{u}^2$.
From Paper~XIV, the edge operator satisfies
$\norm{B_c}_{\op}\leq C'$
uniformly in $c$ on the relevant edge localisation
(the eigenvalues of $B_c$ converge to those of $\cA$,
which are $-a_k = \Theta(k^{2/3})$, bounded on compact index sets;
the operator bound follows from the Airy spectral bound).
Since $H^1(\R_+)\hookrightarrow L^2(\R_+)$ with
$\norm{u}\leq\norm{u}_{H^1}$,
the bound~\eqref{eq:uniformform} holds with
$C = \sup_c\norm{B_c}_{\op} < \infty$.
\end{proof}

\begin{remark}[Role of the uniform bound]
\label{rem:unifrole}
Lemma~\ref{lem:uniformform} is the key ingredient for
Mosco condition~(M1). Without it, the lim inf estimate
for arbitrary weakly convergent sequences cannot be closed
from finite-dimensional form convergence alone.
\end{remark}

\subsection{Mosco convergence}

\begin{definition}[Mosco convergence, cf.~\cite{Mosco1969,Attouch1984}]
\label{def:mosco}
Closed quadratic forms $\mathfrak{q}_c$ on $H$
\emph{Mosco-converge} to $\mathfrak{q}$ if:
\begin{enumerate}[nosep,label=(M\arabic*)]
  \item For every $u_c\rightharpoonup u$ weakly in $H$:
    $\liminf_c\mathfrak{q}_c(u_c,u_c)\geq\mathfrak{q}(u,u)$.
  \item For every $u\in\mathrm{dom}(\mathfrak{q})$:
    $\exists\,u_c\to u$ strongly with
    $\mathfrak{q}_c(u_c,u_c)\to\mathfrak{q}(u,u)$.
\end{enumerate}
Mosco convergence implies strong resolvent convergence
of the associated operators (Attouch--Wets,
cf.~\cite{Attouch1984}).
\end{definition}

\begin{theorem}[Mosco convergence and strong resolvent]
\label{thm:strongres}
Under Assumption~\ref{ass:kernelrate},
$\mathfrak{q}[B_c]$ Mosco-converges to $\mathfrak{q}[\cA]$.
In particular, $(B_c-\zeta)^{-1}\to(\cA-\zeta)^{-1}$ strongly
for all $\zeta\in\C\setminus\R$.
\end{theorem}

\begin{proof}
\textbf{Condition~(M2).}
For $u\in\mathrm{span}\{u_k\}$ (dense in $\mathrm{dom}(\mathfrak{q}[\cA])$
in the form norm), set $u_c$ as in Theorem~\ref{thm:formconv}.
Then $u_c\to u$ strongly in $L^2(\R_+)$
(spectral concentration, Corollary~\ref{cor:concentration}),
and $\mathfrak{q}[B_c](u_c,u_c)\to\mathfrak{q}[\cA](u,u)$
by Theorem~\ref{thm:formconv}.
For general $u\in\mathrm{dom}(\mathfrak{q}[\cA])$,
choose $u^{(K)}\in\mathrm{span}\{u_k\}_{k\leq K}$ with
$\norm{u-u^{(K)}}_{H^1}\leq\varepsilon$.
By Lemma~\ref{lem:uniformform},
$|\mathfrak{q}[B_c](u_c,u_c)-\mathfrak{q}[B_c](u_c^{(K)},u_c^{(K)})|
\leq C\norm{u-u^{(K)}}_{H^1}^2\leq C\varepsilon^2$,
uniformly in $c$. Sending $K\to\infty$ then $\varepsilon\to 0$
gives~(M2) for all $u$.

\textbf{Condition~(M1).}
Let $u_c\rightharpoonup u$ weakly in $L^2(\R_+)$.
Fix $K$ and project: $u_c = P_K u_c + u_c^\perp$
where $P_K$ is the spectral projector of $\cA$ onto
$\mathrm{span}\{u_1,\ldots,u_K\}$.

For the $P_K$-part:
$P_K u_c \to P_K u$ weakly (hence strongly in $K$-dimensional space),
and by Theorem~\ref{thm:formconv}:
$\mathfrak{q}[B_c](P_K u_c,P_K u_c)\to\mathfrak{q}[\cA](P_K u,P_K u)$.

For the remainder $u_c^\perp$:
$\mathfrak{q}[B_c](u_c^\perp,u_c^\perp)
\geq -C\norm{u_c^\perp}^2$
by Lemma~\ref{lem:uniformform}
(the form is bounded below uniformly in $c$).
Since $u_c$ is weakly bounded, $\norm{u_c^\perp}^2$ is bounded;
and as $K\to\infty$, $\mathfrak{q}[\cA](u-P_K u,u-P_K u)\to 0$.

Combining via the off-diagonal control from
Assumption~\ref{ass:kernelrate}
(which prevents cancellation between the $P_K$ and $\perp$ parts
by bounding mixed terms via the $L^2$ kernel rate):
\[
  \liminf_c\mathfrak{q}[B_c](u_c,u_c)
  \geq \mathfrak{q}[\cA](P_K u, P_K u)
  - C'g(c)\norm{u}^2.
\]
Sending $K\to\infty$ and using $g(c)\to 0$ gives
$\liminf_c\mathfrak{q}[B_c](u_c,u_c)\geq\mathfrak{q}[\cA](u,u)$,
i.e.~(M1).
\end{proof}

\begin{remark}[The role of each ingredient in Mosco]
\label{rem:moscoroles}
\begin{itemize}[nosep]
  \item \emph{Theorem~\ref{thm:formconv}} closes~(M2)
    and the $P_K$-part of~(M1).
  \item \emph{Lemma~\ref{lem:uniformform}} gives the
    uniform lower bound needed for the remainder in~(M1).
  \item \emph{Assumption~\ref{ass:kernelrate}} controls
    the off-diagonal mixing between the $P_K$ and $\perp$
    parts in~(M1), preventing spurious cancellation.
\end{itemize}
None of the three can be removed: without the uniform bound,
the remainder could contribute negatively without control;
without the kernel rate, the mixing term is uncontrolled.
\end{remark}

%% ============================================================
\section{The Schaltpunkt: Local Airy Control vs.\ Global Integrability}
\label{sec:switch}
%% ============================================================

\begin{remark}[Operator-theoretic hierarchy]
\label{rem:switch}
The regularity cascade is:
\begin{equation}\label{eq:hierarchy}
  \underbrace{\text{quasimodes}}_{\text{local}}
  \Rightarrow
  \underbrace{\text{form conv.}}_{\text{fin.-dim.}}
  \Rightarrow
  \underbrace{L^2\text{ kernel}}_{\text{switch point}}
  \Rightarrow
  \underbrace{\text{Mosco}\Rightarrow\text{strong resolvent}}_{\text{operator level}}
  \Rightarrow
  \underbrace{\text{trace-norm}}_{\text{Schatten}}
\end{equation}
In operator ideal terms:
$\mathrm{strong} < \mathrm{Mosco/form} < L^2\text{-kernel} < \mathcal{S}_1$.
No arrow reverses; each is witnessed by explicit counterexamples
(Proposition~\ref{prop:nonimpl}).
\end{remark}

\begin{proposition}[Non-implications]
\label{prop:nonimpl}
\begin{enumerate}[nosep,label=(\roman*)]
  \item Finite-dimensional form convergence
    $\not\Rightarrow$ $L^2$ kernel rate.
  \item $L^2$ kernel rate + strong resolvent
    $\not\Rightarrow$ trace-norm convergence.
\end{enumerate}
\end{proposition}

\begin{proof}
(i): The $K\times K$ Gram matrix of $B_c$ on $V_K^{(c)}$
gives no information on the kernel values outside the Airy window
$\xi,\eta\in[0,R_0]$.
(ii): Standard: strong resolvent convergence does not imply Schatten-norm
convergence of resolvents (counterexample:
$A_n=\mathrm{diag}(1/n,2/n,\ldots)\to 0$ strongly,
but $(A_n-i)^{-1}\not\to(-i)^{-1}$ in trace norm).
\end{proof}

%% ============================================================
\section{The Remaining Gap: Strong Resolvent vs.\ Trace-Norm}
\label{sec:gap}
%% ============================================================

\begin{remark}[BR1--BR3 gap is structural]
Strong resolvent convergence does not imply trace-norm convergence.
Closing the gap requires:
\begin{equation}\label{eq:BR3cond}
  \int_{\R_+^2}
  |K_{B_c}(\xi,\eta) - K_{\cA}(\xi,\eta)|^2 w(\xi)w(\eta)\,d\xi\,d\eta
  \to 0,
  \quad w(\xi)=(1+\xi)^{-\beta},
\end{equation}
strictly stronger than Assumption~\ref{ass:kernelrate}.
\end{remark}

\begin{problem}[BR3: weighted kernel convergence]
\label{prob:BR3}
Prove~\eqref{eq:BR3cond} with explicit rate.
\begin{enumerate}[nosep,label=(\roman*)]
  \item Near-diagonal: $|K_{B_c}-K_{\cA}|\leq g(c)h(\xi,\eta)$
    with $h\in L^2(w\otimes w)$.
  \item Off-diagonal: Paley--Wiener decay of $K_{B_c}$
    for $|\xi-\eta|\gg c^{-1/3}$.
\end{enumerate}
\end{problem}

\begin{remark}[Single bottleneck]
\label{rem:bottleneck}
After Papers~XIII--XVI, all difficulty is concentrated in
$\|K_{B_c}-K_{\cA}\|_{L^2(w\otimes w)}$.
Every other step is in place:
scale compatibility (Lemma~\ref{lem:scalecompat}),
spectral identification (Prop.~\ref{prop:specid}),
Davis--Kahan concentration (Cor.~\ref{cor:concentration}),
form convergence (Thm.~\ref{thm:formconv}),
uniform form bound (Lemma~\ref{lem:uniformform}),
Mosco and strong resolvent (Thm.~\ref{thm:strongres}).

One caveat: the ``single bottleneck'' claim is accurate
provided there are no hidden logarithmic losses in the
off-diagonal PSWF tail decay (Problem~\ref{prob:BR3}(ii))
or in the near-diagonal Airy approximation error.
These potential sources of loss are
explicitly isolated in Problem~\ref{prob:BR3}
and constitute the actual content of the BR3 problem.
\end{remark}

\subsection*{Programme position after Papers~XIII--XVI}

\begin{center}
\renewcommand{\arraystretch}{1.35}
\begin{tabular}{p{5.4cm} p{1.5cm} p{2.1cm}}
\toprule
\textbf{Claim} & \textbf{Status} & \textbf{Note} \\
\midrule
Spectral ID, ordered matching (Prop.~\ref{prop:specid})
  & \checkmark & unconditional \\
Scale Compatibility (Lem.~\ref{lem:scalecompat})
  & \checkmark & unconditional \\
Concentration $O(c^{-7/3})$ (Cor.~\ref{cor:concentration})
  & \checkmark & unconditional \\
Form conv.\ on Airy subspaces (Thm.~\ref{thm:formconv})
  & \checkmark & unconditional, $c^{-1/6}$ \\
Uniform form bound (Lem.~\ref{lem:uniformform})
  & \checkmark & unconditional \\
$L^2$ kernel rate (Ass.~\ref{ass:kernelrate})
  & open & decisive switch point \\
Mosco convergence of forms & conditional & Ass.~\ref{ass:kernelrate} \\
Strong resolvent (Thm.~\ref{thm:strongres})
  & conditional & follows from Mosco \\
Weighted kernel conv.\ BR3 (Prob.~\ref{prob:BR3})
  & open & single bottleneck \\
Trace-norm conv.\ (Paper~XIV, BR3)
  & open & needs BR3 \\
(Q5-lower) $N_\varphi\geq 2$ (Paper~XV)
  & \checkmark & unconditional \\
(Q5-upper) $N_\varphi\leq 2$ (Paper~XIV)
  & open & needs BR3 \\
Gap-S (Paper~XIII)
  & open & needs (Q5) \\
\bottomrule
\end{tabular}
\end{center}

\medskip\noindent
\textbf{Logical chain:}
\begin{equation*}
\boxed{
\underbrace{\text{Spec.~ID}+\text{Scale}+\text{Quasimodes}}_{\text{unconditional}}
\xrightarrow{}
\underbrace{\text{BR1}_{\text{fin}}+\text{UFormBd}}_{\text{unconditional}}
\xrightarrow{+\text{Ass.}\ref{ass:kernelrate}}
\underbrace{\text{Mosco}\Rightarrow\text{BR1}_{\text{glob}}}_{\text{switch point}}
\xrightarrow{+\text{Prob.}\ref{prob:BR3}}
\underbrace{(Q5)}_{\text{bottleneck}}
\xrightarrow{}
\text{Gap-S}
}
\end{equation*}

\begin{thebibliography}{99}
\bibitem{paper13} \textit{Paper~XIII}, \texttt{paper13\_gap\_s.tex}, 2026.
\bibitem{paper14} \textit{Paper~XIV}, \texttt{paper14\_airy\_resolvent.tex}, 2026.
\bibitem{paper15} \textit{Paper~XV}, \texttt{paper15\_quasimode.tex}, 2026.
\bibitem{Attouch1984} H.~Attouch, \textit{Variational Convergence
  for Functions and Operators}, Pitman, 1984.
\bibitem{DavisKahan1970} C.~Davis, W.~M.~Kahan,
  The rotation of eigenvectors by a perturbation~III,
  \textit{SIAM J.\ Numer.\ Anal.}\ \textbf{7} (1970), 1--46.
\bibitem{Kato1966} T.~Kato, \textit{Perturbation Theory for Linear Operators},
  Springer, 1966.
\bibitem{Mosco1969} U.~Mosco,
  Convergence of convex sets and of solutions of variational inequalities,
  \textit{Adv.\ Math.}\ \textbf{3} (1969), 510--585.
\bibitem{Osipov2013} A.~Osipov, V.~Rokhlin, H.~Xiao,
  \textit{Prolate Spheroidal Wave Functions of Order Zero}, Springer, 2013.
\bibitem{Olver1997} F.~Olver, \textit{Asymptotics and Special Functions},
  AK Peters, 1997.
\bibitem{ReedSimon1980} M.~Reed, B.~Simon,
  \textit{Methods of Modern Mathematical Physics~I}, Academic Press, 1980.
\bibitem{Simon2005} B.~Simon, \textit{Trace Ideals and Their Applications},
  2nd ed., AMS, 2005.
\end{thebibliography}

\end{document}
